\documentclass[11pt]{article}

\usepackage[margin=1in]{geometry}
\usepackage{amsmath,amssymb}
\usepackage{enumitem}
\usepackage{hyperref}
\usepackage{bm}
\usepackage{booktabs}

\title{Short-Read Isoform Quantification as a GLM}
\author{}
\date{}

\begin{document}
\maketitle

\section*{Purpose}

Common short-read isoform quantification models can be written in a common
generalized linear model (GLM) framework. This formulation also defines the
read-isoform alignment conditional probability matrix used in
quantification-error analysis and K-value calculation.

\section*{Basic Model}

Consider a gene with \(T\) isoforms and relative abundance vector
\[
  \boldsymbol{\theta} = (\theta_1,\theta_2,\ldots,\theta_T)'.
\]
There are \(N\) independent reads \(R_1,\ldots,R_N\), each with common read
length \(L\). For isoform \(t\), let \(l_t\) be the isoform length and
\[
  \tilde{l}_t = l_t - L + 1
\]
be its effective length. Read start sites are assumed to be uniformly
distributed over the isoform of origin.

Under this setup, the probability that a read originates from isoform \(t\) is
\[
  \phi_t =
  \Pr(R_n \in t)
  =
  \frac{\theta_t \tilde{l}_t}
       {\sum_{j=1}^T \theta_j \tilde{l}_j}.
\]

\section*{Read-Based and Region-Based Views}

The read-based likelihood models the probability of observing each read
directly. A read contributes according to the isoforms to which it can map:
\[
  \Pr(\text{observing } R_n)
  =
  \sum_{t=1}^T
  \frac{1}{\tilde{l}_t}
  \mathbf{1}\{R_n \text{ maps to isoform } t\}
  \phi_t.
\]

The region/equivalence class-based likelihood first partitions a gene into disjoint regions such
that reads falling in the same region map to the same set of isoforms (equivalence classes). If region
\(s\) has effective length \(\tilde{u}_s\), the probability that a randomly
selected read falls in region \(s\) is
\[
  p_s =
  \sum_{t=1}^T
  \frac{\tilde{u}_s}{\tilde{l}_t}
  \mathbf{1}\{s \in t\}
  \phi_t.
\]
Region counts can then be modeled with a multinomial distribution, approximated
with independent Poisson variables, or approximated using a normal model for
region read proportions.

\section*{Main Equivalence Result}

In the basic model, read-based and region-based formulations yield equivalent
estimators. The key observation is that reads within the same disjoint region
have identical mapping compatibility and therefore make the same likelihood
contribution. Grouping reads by region turns the read-based likelihood into the
multinomial region likelihood up to a constant. The Poisson approximation is
also equivalent to these formulations, so the read-based, multinomial, and
Poisson likelihoods produce the same maximum-likelihood estimator of
\(\boldsymbol{\theta}\).

The normal model is not written as the same likelihood, but its mean structure
has the same linear form. Thus its estimator approximates the estimators from
the read-based, multinomial, and Poisson models.

\section*{GLM Formulation}

For region read proportions \(c_s\), the normal model uses the identity-link
mean equation
\[
  \mathbb{E}[c_s]
  =
  \sum_{t=1}^T
  \frac{\tilde{u}_s}{\tilde{l}_t}
  \mathbf{1}\{s \in t\}
  \phi_t,
  \qquad 1 \le s \le S.
\]
For region read counts \(n_s\), the multinomial and Poisson formulations use
\[
  \mathbb{E}[n_s]
  =
  N
  \sum_{t=1}^T
  \frac{\tilde{u}_s}{\tilde{l}_t}
  \mathbf{1}\{s \in t\}
  \phi_t,
  \qquad 1 \le s \le S.
\]
Dividing the count equation by \(N\) gives the same mathematical relationship
as the proportion equation. Both are identity-link GLMs relating observed
regional summaries to isoform-origin probabilities.

\section*{Conditional Probability Matrix}

The normal model can be written as a least-squares problem:
\[
  \hat{\boldsymbol{\phi}}
  =
  \arg\min_{\boldsymbol{\phi} \in \mathbb{R}_+^T}
  \left\|
    \mathbf{c} - A\boldsymbol{\phi}
  \right\|_2^2,
\]
where \(\mathbb{R}_+^T\) denotes the nonnegative orthant. Nonnegativity follows
from the interpretation of \(\phi_t\) as an isoform-origin probability or
nonnegative isoform signal. A sum-to-one constraint is optional: if
\(\mathbf{c}\) is normalized and the columns of \(A\) have compatible total
mass, the least-squares loss identifies the scale of
\(\boldsymbol{\phi}\). The fitted vector can then be normalized post hoc if a
probability vector is desired. Here
\[
  \mathbf{c} = (c_1,c_2,\ldots,c_S)'
\]
is the vector of region read proportions and
\[
  A = (a_{st}) \in \mathbb{R}^{S \times T}
\]
is the read-isoform alignment conditional probability matrix. Its entries are
\[
  a_{st}
  =
  \Pr(R_n \in s \mid R_n \in t)
  =
  \frac{\tilde{u}_s}{\tilde{l}_t}
  \mathbf{1}\{s \in t\}.
\]
Equivalently, the estimation problem is the linear system
\[
  \mathbf{c} = A\boldsymbol{\phi}.
\]
This linear system is the basis for defining the K-value.

\section*{Many-Cell Extension and Nuclear-Norm Regularization}

For single-cell data, isoforms can be quantified jointly across \(M\) cells.
The formulation applies either to one gene or to a genome-wide transcript set;
the latter is useful when low-rank structure is intended to capture shared cell
states rather than gene-specific effects. Let
\[
  C = [\mathbf{c}_1,\mathbf{c}_2,\ldots,\mathbf{c}_M]
  \in \mathbb{R}^{S \times M}
\]
be the matrix of region read proportions, with one column per cell, and let
\[
  \Phi = [\boldsymbol{\phi}_1,\boldsymbol{\phi}_2,\ldots,\boldsymbol{\phi}_M]
  \in \mathbb{R}^{T \times M}
\]
be the matrix of cell-specific isoform signals. The entries of \(\Phi\) are
constrained to be nonnegative. If the same read-isoform compatibility matrix
\(A\) is used for all cells, the simultaneous least-squares problem is
\[
  \hat{\Phi}
  =
  \arg\min_{\Phi \ge 0}
  \frac{1}{2}\left\|C - A\Phi\right\|_F^2.
\]
If column-normalized isoform proportions are needed, each fitted column can be
converted to
\[
  \tilde{\boldsymbol{\phi}}_m
  =
  \frac{\hat{\boldsymbol{\phi}}_m}
       {\mathbf{1}'\hat{\boldsymbol{\phi}}_m}
\]
after optimization, for columns with positive total mass.

Single-cell read coverage is sparse, so both \(A\) and \(C\) are typically very
sparse. To share information across cells, one can either penalize the nuclear
norm of the matrix \(\Phi\),
\[
  \hat{\Phi}
  =
  \arg\min_{\Phi \ge 0}
  \frac{1}{2}\left\|C - A\Phi\right\|_F^2
  + \lambda \left\|\Phi\right\|_*.
\]
or constrain it directly:
\[
  \hat{\Phi}
  =
  \arg\min_{\Phi \ge 0,\,\|\Phi\|_* \le \tau}
  \frac{1}{2}\left\|C - A\Phi\right\|_F^2.
\]
The nuclear norm \(\|\Phi\|_*\), the sum of singular values of \(\Phi\), favors
a low-rank structure in which isoform usage across cells is explained by a
small number of shared patterns. 

\subsection*{Regularizing Read-Origin or Molecular Abundance}

The response matrix is normalized columnwise: for each nonempty cell or
pseudobulk, EC counts are divided by their total, so
\(\mathbf{1}'\mathbf{c}_m=1\). UMI depth therefore does not set the scale of a
column or weight its contribution to the least-squares loss.

The linear model can regularize either read-origin signal \(\Phi\) or molecular
abundance signal \(\Theta\), without imposing a simplex constraint. Let
\(D_l=\operatorname{diag}(\tilde l_1,\ldots,\tilde l_T)\). Since
\(\Phi=D_l\Theta\) up to the column scale identified by the normalized
response,
\[
  C \approx A_\phi\Phi
    = A_\phi D_l\Theta
    = A_\theta\Theta.
\]
Thus regularizing \(\Theta\) uses \(A_\theta=A_\phi D_l\) and replaces
\(\lambda\|\Phi\|_*\) by \(\lambda\|\Theta\|_*\). Under the current geometric
approximation, \(A_{\phi,st}=\tilde u_s/\tilde l_t\) and
\(A_{\theta,st}=\tilde u_s\) for compatible isoforms. Fitted columns are
normalized only when producing relative-abundance output; the optimization
itself remains nonnegative but unconstrained in column sum.

\subsection*{Sparse Objective Evaluation}

A scalable implementation should avoid forming the dense residual
\(C-A\Phi\). The squared loss can be expanded as
\[
  \left\|C - A\Phi\right\|_F^2
  =
  \left\|C\right\|_F^2
  - 2\,\operatorname{tr}\!\left(\Phi' A' C\right)
  + \operatorname{tr}\!\left(\Phi' A' A \Phi\right).
\]
Thus the data enter through
\[
  B = A'C,
  \qquad
  G = A'A.
\]
For a small problem, both can be computed with sparse matrix multiplication.
For a genome-wide single-cell matrix, however, \(B\) is generally dense even
when \(A\) and \(C\) are sparse, so it must not be materialized. The gradient
of the smooth loss is
\[
  \nabla_{\Phi}
  \frac{1}{2}\left\|C-A\Phi\right\|_F^2
  =
  G\Phi - B,
\]
which requires multiplication by the isoform-by-isoform matrix \(G\), rather
than by a large dense region-by-cell residual. When
\(\Phi=UV'\), with \(U\in\mathbb{R}^{T\times r}\) and
\(V\in\mathbb{R}^{M\times r}\), the required products can be evaluated over
cell blocks \(C_b\) and \(V_b\):
\[
  A'C_bV_b,\qquad C_b' AU,\qquad U'GU,\qquad V_b'V_b.
\]
Thus the implementation streams sparse cell-by-equivalence-class blocks and
keeps only the \(T\times r\) and \(M\times r\) factors in memory. This avoids
forming \(B\), the residual, or the full \(T\times M\) abundance matrix during
optimization.

\subsection*{Optimization Strategies}

There are three practical routes, plus a factorized ADMM variant for the
large-scale setting.

\begin{itemize}[leftmargin=*]
  \item \textbf{ADMM with a nuclear-norm split.} Introduce a second matrix
  \(Z\) and solve
  \[
    \min_{\Phi \ge 0,\,Z}
    \frac{1}{2}\left\|C-A\Phi\right\|_F^2
    + \lambda\|Z\|_*
    \quad\text{subject to}\quad
    \Phi = Z.
  \]
  The \(\Phi\)-update is a sparse nonnegative quadratic problem. The
  corresponding unconstrained normal equation has the form
  \[
    (A'A + \rho I)\Phi = A'C + \rho(Z-U),
  \]
  where \(U\) is the scaled dual variable and \(\rho\) is the ADMM penalty.
  Therefore the update can use the sparse products \(B=A'C\) and \(G=A'A\)
  instead of forming the dense residual \(C-A\Phi\). If \(T\) is modest for a
  gene, factorize \(G+\rho I\) once and reuse it for all cells. If \(T\) is
  large, solve the update iteratively using products with sparse \(A\) and
  \(A'\), or with \(G\) if \(G\) is still sparse enough. The nonnegativity
  constraint can be handled by projected conjugate gradient, coordinate
  descent, active-set nonnegative least squares, or an additional proximal
  projection step.

  The \(Z\)-update is singular-value soft-thresholding:
  \[
    Z \leftarrow
    \operatorname{SVT}_{\lambda/\rho}(\Phi + U).
  \]
  This step acts on the isoform-by-cell matrix and does not involve \(A\) or
  \(C\). Thus ADMM can exploit sparsity in the data-fitting step, but the
  nuclear-norm proximal step may become the bottleneck when \(M\) is very
  large. The implemented convex ADMM solver is therefore a bounded reference
  method: it stores \(\Phi\), \(Z\), and the dual variable densely, uses a
  projected-gradient inner solve for the nonnegative update, and rejects
  problems beyond a configured dense-entry limit.

  A separate \emph{factorized ADMM} mode replaces the dense state by
  \(U,V\) and nonnegative split copies of these factors. It alternates
  factor updates with proximal nonnegative projections and stores
  \(O(r(T+M))\) state. This has the same scalability advantages as the
  factorized formulation below, but it is non-convex and is not equivalent to
  the bounded convex ADMM reference.

  Both ADMM implementations can adapt \(\rho\) by residual balancing. Every
  configured interval, let \(r\) be the primal split residual and \(s\) the
  dual residual. If \(\|r\|>\eta\|s\|\), multiply \(\rho\) by \(\kappa\); if
  \(\|s\|>\eta\|r\|\), divide it by \(\kappa\). The scaled dual variable is
  multiplied by the old-to-new \(\rho\) ratio so that its unscaled value is
  unchanged. The implementation defaults to \(\eta=10\), \(\kappa=2\), and an
  update every ten iterations, and records the complete \(\rho\) history.

  \item \textbf{Frank-Wolfe with penalized nonnegativity.} Direct
  Frank--Wolfe over
  \[
    \{\Phi:\Phi\geq 0,\ \|\Phi\|_*\leq\tau\}
  \]
  does not have the usual singular-vector linear oracle. In particular,
  replacing the signed descent matrix by its entrywise positive part is not an
  exact oracle: a nonnegative rank-one outer product can span both positive and
  negative entries of the original matrix. The legacy \texttt{frank\_wolfe}
  implementation makes this approximation and its reported gap is therefore
  not a valid certificate for the intersection problem.

  The scalable corrected variant instead solves
  \[
    \min_{\|\Phi\|_*\le\tau}
    \frac{1}{2}\|C-A\Phi\|_F^2
    +\frac{\mu}{2}\|\Phi_-\|_F^2,
    \qquad \Phi_-=\min(\Phi,0).
  \]
  Its negative gradient is
  \[
    Q_k=B-G\Phi_k-\mu(\Phi_k)_-.
  \]
  The ordinary nuclear-norm ball has the signed rank-one oracle
  \[
    S_k=\tau u_kv_k',
  \]
  where \(u_k,v_k\) are leading singular vectors of \(Q_k\). The implementation
  obtains them with matrix-free power iterations on the signed operator rather
  than materializing \(B\), \(Q_k\), or \(\Phi_k\). Sparse products evaluate
  the least-squares terms, while transcript-by-cell blocks are reconstructed
  only to apply \((\Phi_k)_-\).

  The update
  \[
    \Phi_{k+1}=(1-\gamma_k)\Phi_k+\gamma_kS_k
  \]
  preserves the nuclear-norm constraint. A curvature-bounded step uses
  \(\|A(S_k-\Phi_k)\|_F^2+\mu\|S_k-\Phi_k\|_F^2\). The penalty coefficient is
  specified relative to the estimated \(\|A\|_2^2\), making its scale portable
  across design matrices. Finite \(\mu\) gives approximate rather than exact
  nonnegativity, so manifests report negative Frobenius and absolute-mass
  fractions. The power-iteration candidate gap is a stopping diagnostic but is
  explicitly not labeled an exact dual certificate.

  \item \textbf{Low-rank factorization.} Use the variational identity
  \[
    \|\Phi\|_*
    =
    \min_{\Phi=UV'}
    \frac{1}{2}\left(\|U\|_F^2 + \|V\|_F^2\right)
  \]
  and optimize a fixed-rank approximation \(\Phi \approx UV'\), where
  \(U \in \mathbb{R}^{T \times r}\), \(V \in \mathbb{R}^{M \times r}\), and
  \(r \ll \min(T,M)\). This avoids full singular-value decompositions and
  stores only \(O(r(T+M))\) parameters. The loss can be evaluated using \(B\)
  and \(G\):
  \[
    \left\|C-AUV'\right\|_F^2
    =
    \|C\|_F^2
    -2\,\operatorname{tr}\!\left(VU'B\right)
    +\operatorname{tr}\!\left((U'GU)(V'V)\right).
  \]
  The implementation uses nonnegative factors and projected alternating
  updates. Gradients for the shared transcript factor are accumulated over
  sparse cell blocks and normalized by the number of nonempty cells; cell
  factors are updated blockwise. Column normalization is applied after fitting
  when isoform proportions are required.
\end{itemize}

The factorized and Frank-Wolfe approaches are the scalable choices for many
cells. Both keep \(\Phi\) in a compact low-rank representation and stream
sparse blocks rather than precomputing \(B=A'C\). ADMM is useful as a direct
convex reference for smaller matrices, but its repeated singular-value
thresholding and dense primal state make it unsuitable for a million-cell
matrix.

\subsection*{Software Interface and Output}

The single-cell command exposes the methods as
\texttt{--quant\_method admm}, \texttt{admm\_factorized},
\texttt{frank\_wolfe}, \texttt{frank\_wolfe\_penalized}, and
\texttt{factorized}. The scalable methods use a
PyTorch backend and select CUDA when available; NumPy/SciPy remains responsible
for sparse input and compatibility-matrix construction. The rank, cell block
size, convergence tolerance, regularization strength, ADMM penalty, and
Frank--Wolfe radius are configurable. The default low-rank capacity is 64, and
the default Frank--Wolfe radius is \(\sqrt{M}\). The
\texttt{--regularization\_target} option selects \texttt{phi} or
\texttt{theta}; it changes both the coefficient matrix receiving the penalty
and the corresponding design matrix.

In \texttt{single\_cell} mode, fitting operates directly on the raw
cell-by-equivalence-class matrix rather than creating pseudobulks. The output
consists of compact low-rank factors, barcode and transcript order, and a
manifest of solver diagnostics. Thresholded sparse transcript-by-cell chunks
are optional because even sparse reconstruction can be large when many fitted
values exceed the threshold. Selected cell blocks can instead be reconstructed
from the saved factors. The downstream intron clustering workflow remains a
pseudobulk operation and is not invoked in this mode.

Because each response column has unit mass, both random factors are initialized
at scale \(1/\sqrt{rT}\). This keeps the initial fitted mass at order one;
scaling only by rank initializes the fitted mass at order \(T\) and can drive
a projected method to the all-zero solution on its first update.
Projected factor updates use curvature-bounded step sizes rather than one fixed
step across datasets. The cell-factor bound is computed from the rank-sized
matrix \(U'A'AU\). The transcript-factor bound combines \(V'V\) with a sparse
power-iteration estimate of \(\lVert A\rVert_2^2\); the ADMM bounds also include
the split penalty \(\rho\). This prevents an update from crossing the
nonnegative orthant and collapsing an entire factor to zero when the design
curvature is large.

The factorized solvers stop only after a configured minimum iteration count
and several consecutive relative-objective changes below tolerance.
Factorized ADMM additionally requires its relative primal and dual split
residuals to satisfy the tolerance. Penalized Frank--Wolfe requires both
relative objective change and its approximate candidate gap to remain below
tolerance; this candidate gap is not presented as a certificate.
Its atom capacity is configured separately from the fixed rank used by the
factorized methods, and additional power iterations can improve the approximate
linear oracle. Every manifest records the stopping history, convergence reason,
and whether the maximum iteration or atom limit was reached.

For comparison with a standard single-cell expression analysis, downstream
label scoring does not apply a nonlinear transform to the arbitrary factors.
Instead, transcript loadings are summed by gene so that fitted gene abundance
remains \(VU_{\mathrm{gene}}'\). For each cell this abundance is normalized to
a total of 10,000 and transformed with \(\log(1+x)\). Per-fit gene moments are
computed in streamed cell blocks, the 2,000 most variable genes are selected,
and incremental PCA is fit and transformed in two further streamed passes.
Grouped classification, k-means, ARI, NMI, and silhouette are all computed in
the resulting 30-dimensional PCA representation. This preserves the intended
nonlinear expression preprocessing without materializing a dense genome-wide
cell-by-gene matrix.

\subsection*{Choosing the Regularization Level}

The implemented cross-validation partitions molecule counts rather than rows
of \(A\). Each observed integer EC count is randomly assigned to one of \(Q\)
folds. For fold \(q\), the remaining counts are normalized per cell and used to
fit \(\Phi\); the held-out counts are normalized independently and scored with
the same design matrix:
\[
  \operatorname{CV}_q(h)
  = \frac{1}{M_q}
    \left\|C_{\mathrm{val}}^{(q)}
    -A\hat\Phi_h^{(q)}\right\|_F^2.
\]
This tests prediction of independent molecules from the same cells while
retaining every compatibility pattern in the design and avoiding a separate
procedure for unseen cell factors. The current full-data workflow uses three
folds on a reproducible cell subset, chooses the multiplier with minimum mean
held-out loss, and then refits all cells.

Candidate values are dimensionless so that a subset-selected multiplier can be
transferred to the complete dataset. For the penalized objective, define
\[
  \lambda_{\max}=\|A'C\|_2
\]
and tune \(\lambda/\lambda_{\max}\). For constrained Frank--Wolfe, let \(u,v\)
be the leading singular vectors of \(B=A'C\), with singular value \(\sigma\),
and define the rank-one line-search reference
\[
  \tau_1=\frac{\sigma}{\|Au\|_2^2}.
\]
The search tunes \(\tau/\tau_1\). Both scales are estimated by streamed power
iterations without materializing \(B\). After selecting a multiplier on the
subset, the corresponding scale is recomputed using all cells and their product
is passed as the fixed \(\lambda\) or \(\tau\) for the final fit.
The factorized implementation averages the complete left-factor block gradient,
including its regularization and augmented-Lagrangian terms, by the number of
nonempty cells. This preserves the relative scale of the objective terms when a
multiplier selected on a subset is transferred to all cells. An optimum on an
open boundary of the candidate grid triggers automatic geometric expansion.
Only the new endpoint is fit in each round and its fold losses are merged with
the existing candidates. Expansion stops at an interior optimum, at zero or
\(\lambda/\lambda_{\max}=1\), or at a configured expansion limit; the full fit
is rejected only in the last case.

For the penalized Frank--Wolfe solver, model selection uses the
one-standard-error rule. Let \(\hat h\) minimize mean validation loss and let
\(s_{\hat h}\) be its fold standard error. Among converged candidates satisfying
\(\overline{L}(h)\leq\overline{L}(\hat h)+s_{\hat h}\), the smallest nuclear
radius is selected. This favors the most regularized statistically
indistinguishable fit and prevents negligible endpoint improvements from
driving indefinite grid expansion. The penalized oracle's candidate gap is
diagnostic rather than a convergence certificate, so stopping uses patient
relative objective change.

Statistical equivalence in reconstruction loss does not make a degenerate
representation useful. Before applying the one-standard-error rule, each
candidate is therefore screened after rectification, per-cell library
normalization, and a log-one-plus transform on a fixed set of high-loading
features. Every fold must retain the configured active-cell fraction and exceed
a minimum relative between-cell profile variance. This scale-invariant check
excludes zero solutions and rank-one factors that vary only by library size.

Among the converged, nondegenerate candidates inside the reconstruction
one-standard-error set, define \(v(h)\) as mean normalized profile variance.
The implementation retains candidates satisfying
\[
  v(h) \geq \gamma \max_{h'} v(h'),
\]
with \(\gamma=0.9\) by default, and chooses the smallest retained nuclear
radius. This preserves most observable cell-to-cell structure while retaining
the strongest regularization compatible with both structure and held-out count
prediction. Cell-type labels remain an external assessment rather than a
hyperparameter input.

\section*{Connecting to equivalence-class-based inference}

The preceding derivation used geometric regions with entries
\(\tilde{u}_s/\tilde{l}_t\). Salmon's range-factorized equivalence classes can
be used in the same linear system by replacing those geometric entries with
Salmon's conditional assignment weights.

Let \(K\) be the number of equivalence classes. Running Salmon with
\texttt{--dumpEqWeights} provides the two quantities needed to instantiate the
model:
\begin{itemize}[leftmargin=*]
  \item \texttt{eq\_classes.txt} gives the count \(c_k\) for each class and the
  conditional weights
  \[
    p_{kt} = \Pr(R_n \in k \mid R_n \in t),
  \]
  with \(p_{kt}=0\) when transcript \(t\) is incompatible with class \(k\).
  These weights define the design matrix \(A \in \mathbb{R}^{K \times T}\),
  \(a_{kt}=p_{kt}\).
  \item \texttt{quant.sf} gives the effective length vector
  \(\tilde{\bm{l}}=(\tilde{l}_1,\ldots,\tilde{l}_T)'\), including Salmon's
  fragment-length and bias corrections.
\end{itemize}

With this substitution, the isoform-origin formulation is exactly the one above:
counts use \(\mathbb{E}[\bm{c}]=N A\bm{\phi}\), or equivalently
\(\bm{c}=A\bm{\phi}^*\) for the unnormalized signal
\(\bm{\phi}^*=N\bm{\phi}\). The same nonnegative least-squares and many-cell
extensions therefore apply without explicitly computing geometric region
lengths \(\tilde{u}_s\).

The current single-cell implementation constructs \(A\) from EC membership
when only the standard alevin-fry EC dump is present: compatible entries are
weighted by \(1/\tilde{l}_t\), using transcript effective-length proxies from
the reference, and the region factor is set to one. The patched weighted RAD
and alevin-fry route additionally writes per-UMI alignment-likelihood vectors
to \texttt{gene\_eqclass\_probs.tsv.gz}. These vectors are computed from
alignment score and transcript-end position and are used inside alevin-fry's
EM update. They vary by cell and UMI; they are not yet consumed by the GLM,
which currently assumes one common \(A\) across cells. Using them exactly
requires a cell-specific design operator rather than replacing \(A\) by a
single averaged matrix.

For a direct fixed-design comparison, the implementation also supports a
globally weighted approximation. For each global EC, its per-UMI likelihood
vectors are averaged over cells and UMIs; ECs without dumped vectors receive
uniform weights over their compatible transcripts. The resulting sparse matrix
is column-normalized over ECs for each transcript and cached. The binary
comparator column-normalizes the binary EC membership matrix in the same way.
This comparison isolates the effect of within-EC alignment evidence while
retaining one shared design matrix. It is an empirical approximation to the
cell-specific likelihoods, not an annotation-only conditional model.

If inference is desired directly on the relative abundance scale
\(\boldsymbol{\theta}\), absorb transcript effective lengths into the design:
\[
  B = A\operatorname{diag}(\tilde{\bm{l}}),
  \qquad
  b_{kt}=p_{kt}\tilde{l}_t.
\]
Solving the corresponding nonnegative system
\(\bm{c}=B\boldsymbol{\theta}^*\) estimates an unnormalized abundance signal;
normalizing its entries to sum to one recovers relative abundances.


\end{document}
